\documentclass[10pt]{report}

\usepackage[spanish, es-tabla]{babel}
\usepackage[utf8]{inputenc}
\usepackage{lmodern}
\usepackage[T1]{fontenc}
\usepackage{setspace}
\usepackage{geometry}
\usepackage{titlesec}
\usepackage{caption}
\usepackage{csquotes}
\usepackage{siunitx}
\usepackage{fancyhdr}
\usepackage[backend=biber, style=apa]{biblatex}
\addbibresource{referencias.bib}

\geometry{letterpaper, margin=2.54cm}

\raggedright
\doublespacing
\unaccentedoperators

\titleformat{\chapter}[block]{\normalfont\bfseries\filcenter\Large}{\thechapter}{1em}{}
\titlespacing*{\chapter}{0pt}{-20pt}{20pt}
\titleformat{\section}[block]{\normalfont\bfseries\large}{\thesection}{1em}{}
\titleformat{\subsection}[block]{\normalfont\bfseries\normalsize}{\thesubsection}{1em}{}

\setlength{\parindent}{1.27cm}

\pagestyle{fancy}
\fancyhf{}
\renewcommand{\headrulewidth}{0pt}
\fancyhead[R]{\thepage}
\fancypagestyle{plain}{
  \fancyhf{}
  \fancyhead[R]{\thepage}
  \renewcommand{\headrulewidth}{0pt}
}

% === ALIAS DE ESPECIE (personalizar) ===
\newcommand{\cmag}{\textit{C. magdalenae}}
\newcommand{\ea}{\textit{Ea}}

\begin{document}

% === PORTADA 1 ===
\begin{titlepage}
    \centering
    {\bfseries\large TÍTULO DEL TRABAJO \par}
    \vfill
    {\bfseries\large NOMBRE DEL AUTOR \par}
    \vfill
    {\bfseries\large UNIVERSIDAD \\
     FACULTAD \\
     PROGRAMA \\
     Ciudad, País \\
     AÑO \par}
\end{titlepage}

% === PORTADA 2 ===
\begin{titlepage}
    \centering
    {\bfseries\large TÍTULO DEL TRABAJO \par}
    \vfill
    {\bfseries\large NOMBRE DEL AUTOR \par}
    \vfill
    {\bfseries\large Trabajo de grado para optar al título de TÍTULO}
    \vfill
    {\bfseries\large DIRECTOR: \\
    Título. Nombre del Director \\
    Cargo, Institución \par}
    \vfill
    {\bfseries\large UNIVERSIDAD \\
     FACULTAD \\
     PROGRAMA \\
     Ciudad, País \\
     AÑO \par}
\end{titlepage}

% === NOTA DE ACEPTACIÓN ===
\newpage
\begin{titlepage}
\begin{flushright}
    {\bfseries NOTA DE ACEPTACIÓN \par}
    \vspace{0.5cm}
    \rule{8cm}{0.4pt} \\
    \vspace{0.2cm}
    \rule{8cm}{0.4pt} \\
    \vspace{0.2cm}
    \rule{8cm}{0.4pt} \\
    \vspace{0.2cm}
    \rule{8cm}{0.4pt}
\end{flushright}
\vfill
\begin{flushright}
    {\bfseries NOMBRE DEL DIRECTOR \par}
    Institución. Director
\end{flushright}
\vfill
\begin{flushright}
    \rule{8cm}{0.4pt} \\
    Firma del Jurado
    
    \vspace{2cm}
    \rule{8cm}{0.4pt} \\
    Firma del Jurado
\end{flushright}
\vfill
\begin{flushright}
    {\bfseries Ciudad – País, AÑO}
\end{flushright}
\end{titlepage}

% === DEDICATORIA ===
\begin{titlepage}
\begin{center}
    {\bfseries\large DEDICATORIA \par}
\end{center}
\raggedleft
Texto de dedicatoria.
\end{titlepage}

% === AGRADECIMIENTOS ===
\begin{titlepage}
\begin{center}
    {\bfseries\large AGRADECIMIENTOS \par}
\end{center}
Texto de agradecimientos.
\end{titlepage}

% === TABLAS DE CONTENIDO ===
\tableofcontents
\listoffigures
\renewcommand{\listtablename}{Índice de tablas}
\listoftables

% === RESUMEN ===
\chapter*{RESUMEN}

Texto del resumen en español.\\

\textbf{Palabras clave:} palabra1, palabra2, palabra3.

% === ABSTRACT ===
\chapter*{ABSTRACT}

Abstract text in english.\\

\textbf{Keywords:} keyword1, keyword2, keyword3.

% === CAPÍTULOS ===

\chapter{INTRODUCCIÓN}

Texto de introducción.

\chapter{OBJETIVOS}

\section{Objetivo general}

Texto del objetivo general.

\section{Objetivos específicos}

\begin{enumerate}
    \item Objetivo específico 1.
    \item Objetivo específico 2.
    \item Objetivo específico 3.
\end{enumerate}

\chapter{MARCO REFERENCIAL}

\section{Antecedentes}

Texto de antecedentes.

\section{Marco teórico}

\subsection{Subtema 1}

Texto.

\subsection{Subtema 2}

Texto.

\chapter{DISEÑO METODOLÓGICO}

\section{Área de estudio}

\section{Fase de campo}

\section{Fase de laboratorio}

\section{Análisis de la información}

\chapter{RESULTADOS}

\chapter{DISCUSIÓN}

\chapter{CONCLUSIONES}

\chapter{REFERENCIAS}

\printbibliography[title={REFERENCIAS}, heading=bibnumbered]

\chapter{ANEXOS}

\end{document}
